\documentclass[english, parskip=full]{scrartcl}
\usepackage[utf8]{inputenc}  % Kodierung der Datei
\usepackage[english]{babel}  % Sprache auf Deutsch 
\usepackage{color}
\usepackage{graphicx, gensymb}
\usepackage{amsfonts, cancel, amssymb}
\usepackage{amsmath, amsthm, array, esint}
\usepackage{booktabs, multirow}  % schönere Tabellen
\usepackage{caption}  % Anpassung der Captions
\usepackage{scrlayer-scrpage}    % Manipulation des Seitenstils
%\pagestyle{scrheadings}  % Stil für die Seite setzen
%\automark{section}  % Abschnittsnamen als Seitenbeschriftung 
%\setheadsepline{.5pt}    % Kopzeile durch Linie abtrennen
\usepackage[hidelinks]{hyperref}  % Links und weitere PDF-Features
\usepackage{typearea, tikz, pgffor, verbatim}
\usetikzlibrary{calc, quotes}
\usepackage[cache=false, outputdir=build]{minted}
\usepackage{placeins} % provides \FloatBarrier

\definecolor{bg}{rgb}{0.9,0.9,0.9}
\newcommand\numberthis{\addtocounter{equation}{1}\tag{\theequation}}
\newcommand{\bra}{}
\newcommand{\kom}{}
\newcommand{\brak}{}
\newcommand{\braket}{}
\newcommand{\intx}{}
\newcommand{\intr}{}
\newcommand{\kla}{}
\newcommand{\klam}{}
\newcommand{\klammer}{}
\newcommand{\oper}{}
\newcommand{\tecxt}{}
\newcommand{\Bra}{}
\newcommand{\Ket}{}
\renewcommand{\i}{\text{i}}
\newcommand{\e}{\text{e}}
\newcommand*{\Fourier}{\textsc{Fourier}\ }
\newcommand*{\F}{\mathcal{F}}
\newcommand*{\sinc}{\text{sinc\,}}
\newcommand{\conv}{\mathop{\scalebox{3.5}{\raisebox{-0.38ex}{\makebox[0.5\width][c]{$\ast$}}}}\limits}

\newcommand*{\titel}{07 Spline-Interpolation und Newton-Verfahren}
\newcommand*{\autor}{Joachim Schwardt und Julian Fleck}

\hypersetup{pdfauthor={\autor}, pdftitle={\titel}}

%\titlehead{titlehead}
\subject{Numerik I - Aufgaben}
\title{\titel}
%\author{\autor}
\author{\begin{tabular}{l|l}
Joachim Schwardt & 4768711 \\ 
Julian Fleck & 4759587\\ 
\end{tabular}}
%\author{\begin{tabular}{|l|l|}
%\hline
%Joachim Schwardt & 4768711 \\ \hline
%Julian Fleck & 4759587\\ \hline
%\end{tabular}}
\date{\today}

\begin{document}

%\begin{titlepage}
\maketitle  % Titel setzen
%\tableofcontents  % Inhaltsverzeichnis setzen
%\end{titlepage}

\section*{Aufgabe 1}
\subsection*{a)}
Gegeben sind zwei Polynome $p$ und $q$ vom Grad kleiner gleich 3. 
Für ein $x_0\in\mathbb{R}$ gelte
\begin{align}
p(x_0) &= q(x_0), &&& p'(x_0) &= q'(x_0), &&& p''(x_0) &= q''(x_0), &&& p'''(x_0) &= q'''(x_0).
\end{align}
Wir können die Polynome allgemein als
\begin{align}
p(x) &= \sum_{k=0}^3 p_k x^k &&\tecxt\text{und} & q(x) &= \sum_{k=0}^3 q_k x^k 
\end{align}
schreiben. 
%Für $n\in\{0,1,2,3\}$ lauten die Gleichungen dann 
%\begin{align}
%p^{(n)}(x_0) &= \sum_{k=0}^{3-n} \frac{(k+n)!}{k!} p_{k+n}x_0^{k} = \sum_{k=0}^{3-n} \frac{(k+n)!}{k!} q_{k+n}x_0^{k} = q^{(n)}(x_0).
%\end{align}
Aus der letzten Gleichung folgt dann direkt $p_3 = q_3$.
Die vorletzte Gleichung ist
\begin{align}
3\cdot 2 p_3x_0 + 2\cdot 1 p_2 &= 3\cdot 2 q_3x_0 + 2\cdot 1 q_2,
\end{align}
wobei man durch Einsetzen der vorherigen Gleichung direkt $p_2 = q_2$ erhält.
Analog findet man dann auch $p_1=q_1$ und $p_0=q_0$.
Damit sind die Polynome $p$ und $q$ identisch, weshalb natürlich dann auch $p(x) = q(x)$ für alle $x\in\mathbb{R}$ gilt.

\subsection*{b)}
Für eine beliebige Funktion ist der kubische Spline typischerweise nur zweimal stetig differenzierbar, hat also Glattheit $\ell=2$. 
Hier existiert dann kein entsprechender Spline aus $S^3_3(\Delta)$.
Formal können wir die Menge natürlich trotzdem hinschreiben, also
\begin{align}
S^3_3 &= \klammer\left\{ s\in C^3(\Delta)\,\middle\vert\, \forall k=0,\dots,N-1\forall x\in [x_k,x_{k+1}): s(x) = \sum_{n=0}^3 a_{k,n} x^n \right\}.
\end{align}
Dabei sind die Anschlussbedingungen bereits durch $s\in C^3$ gefordert.
%Jetzt müssen wir aber noch die Stetigkeit der Polynome und ihrer Ableitungen an den Stützstellen $x_k$ fordern.
Es gibt $4\cdot N$ Parameter $a_{k,n}$, allerdings haben wir auch $4\cdot (N-1)$ Gleichungen für die Anschlussbedingungen.
Das geht nur gut, solange wir höchstens $4$ Parameter ändern müssen, um ein Spline an eine Funktion $f$ anzupassen.
Das werden typischerweise die 4 Koeffizienten auf dem ersten Intervall $[a,x_1)$ sein.
Die Funktion muss also selbst bereits ein Polynom dritten Grades (oder kleiner) sein, womit Spline und Polynom übereinstimmen.
Daher gilt also
\begin{align}
S_3^3 &\equiv \Pi_3.
\end{align}
%Überhaupt hängt doch $S^\ell_m(\Delta)$ auch von der Funktion ab, die interpoliert wird.. ?


\newpage
\section*{Aufgabe 2}
Gegeben ist eine Zerlegung $\{x_k\}$ des Intervalls $I=[a,b]$ mit  $a=x_0<\dots<x_n=b$ und eine auf $I$ zweimal stetig differenzierbare Funktion $f$.
Sei $s$ der zugehörige lineare $C^0$-Spline.
Auf dem Intervall $I_k := [x_k,x_{k+1})$ gilt für alle $x\in I_k$ also
\begin{align}
s(x) &= f(x_k) + (x - x_k)\frac{ f(x_{k+1}) - f(x_k)}{x_{k+1} - x_k} = f(x_k) + (x-x_k) f'(\xi_k),
\end{align}
wobei es gemäß des Mittelwertsatzes ein $\xi_k\in I_k$ geben muss, das die zweite Gleichung erfüllt.
Der Satz über Interpolationsfehler garantiert, dass es ein $\chi_{k,x}\in I_k$ gibt, sodass 
\begin{align}
f(x) - s(x) &= \frac{1}{2}f''(\chi_{k,x}) (x-x_k)(x-x_{k+1}).
\end{align}
Für den Fehler folgt damit natürlich direkt
\begin{align}
\max_{x\in I_k} |f(x) - s(x)| &\le \frac{h_k^2}{2}\max_{x\in I_k} f''(x),
\end{align}
wobei $h_k := x_{k+1} - x_k$.
Insbesondere gilt $|x-x_k|\le h_k$ und $|x-x_{k+1}|\le h_k$, was wir hier ausgenutzt haben.
Eventuell könnte man den numerischen Vorfaktor noch genauer bestimmen, da die Gleichheit nicht für beide Differenzen gleichzeitig auftreten kann.
Für $I$ gilt dann entsprechend
\begin{align}
\max_{x\in I} |f(x) - s(x)| &\le \frac{h_\tecxt\text{max}^2}{2}\max_{x\in I} f''(x),
\end{align}
wobei $h_\tecxt\text{max} = \max_k h_k$.

%Insbesondere gilt nach Taylor aber
%\begin{align}
%f(x) &= f(x_k) + (x-x_k)f'(x_k) + \frac{1}{2}f''(\chi_k)(x-x_k)^2
%\end{align}
%für ein $\chi_{k,x}\in I_k$.
%Daraus folgt dann
%\begin{align}
%\max_{x\in I_k} |f(x) - s(x)| &\le |f'(\xi_k) - f'(x_k)| \max_{x\in I_k} | x-x_k| + \frac{1}{2}\max_{x\in I_k} |x-x_k|^2 f''(\chi_{k,x}) \le \\
%&\le \frac{h_k^2}{2} \max_{x\in I_k} f''(x) + h_k |f'(\xi_k) - f'(x_k)|
%\end{align}
%Insbesondere gilt nach Taylor aber
%\begin{align}
%f(x) &= f(x_k) + (x-x_k) f'(\chi_{k,x})
%\end{align}
%für ein $\chi_{k,x}\in I_k$ (also für jedes $x$ ein anderes $\chi_{k,x}$).
%Sei $h_k := x_{k+1} - x_k$, wobei $|\chi_{k,x} - \xi_k| \le h_k$ gilt (beide Punkt liegen in $I_k$, ihre Differenz muss also kleiner als die Länge des Intervalls sein).
%Daraus folgt dann
%\begin{align}
%\max_{x\in I_k} |f(x) - s(x)| &= \max_{x\in I_k} |x-x_k| \left\vert f'(\chi_{k,x}) - f'(\xi_k) \right\vert = \\
%&= h_k \max_{x\in I_k} \left\vert \intx\int_{\xi_k}^{\chi_{k,x}} f''(t)\, \mathrm{d}t \right\vert \le \\
%&\le h_k \max_{x\in I_k} |\chi_{k,x} - \xi_k| \max_{x\in I_k} f''(x) \le \\
%&\le h_k^2 \max_{x\in I_k} f''(x).
%\end{align}


\newpage
\section*{Aufgabe 3}
Für das Integral einer Funktion $f$ auf dem Intervall $[-1,1]$ 
\begin{align}
I(f) &= \intx\int_{-1}^{1} f(x)\, \mathrm{d}x
\end{align}
möchten wir eine Quadraturformel 
\begin{align}
Q(f) &= a_0f(x_0) + a_1 f(x_1)
\end{align}
mit $x_1 > x_0$ finden, die beliebige Polynome bis zu einem möglichst hohen Grades exakt berechnet.
Wir nehmen zunächst an, dass wir kubische Polynome noch berechnen können.
Sei also $f(x) = b+cx+dx^2+ex^3$, und somit
\begin{align}
I(f) &= bx + \frac{cx^2}{2} + \frac{dx^3}{3} + \frac{ex^4}{4} \Big\vert_{-1}^1 = 2b + \frac{2d}{3}.
\end{align}
Andererseits ist
\begin{align}
Q(f) &= a_0 \kla\left( b+cx_0+dx_0^2+ex_0^3 \right) + a_1 \kla\left( b+cx_1+dx_1^2+ex_1^3 \right) = \\
&= b\kla\left( a_0 + a_1 \right) + c \kla\left( a_0x_0 + a_1x_1 \right) + d\kla\left( a_0x_0^2 + a_1x_1^2 \right) + e\kla\left( a_0x_0^3 + a_1x_1^3 \right).
\end{align}
Da die Koeffizienten $b,c,d$ beliebig sind, muss entsprechend dann
\begin{align}
2 &= a_0 + a_1,\\
0 &= a_0x_0 + a_1x_1,\\
\frac{2}{3} &= a_0x_0^2 + a_1x_1^2,\\
0 &= a_0x_0^3 + a_1x_1^3
\end{align}
gelten. 
Aus der ersten Gleichung folgt $a_1 = 2 - a_0$.
Einsetzen in die zweite Gleichung ergibt
\begin{align}
0 &= a_0x_0 + (2-a_0)x_1 && \Rightarrow & a_0 &= \frac{2x_1}{x_1 - x_0},
\end{align}
womit $a_1 = -\frac{2x_0}{x_1 - x_0}$ ist.
Einsetzen in die dritte Gleichung ergibt dann
\begin{align}
\frac{2}{3} &= \frac{2x_1}{x_1 - x_0}x_0^2 - \frac{2x_0}{x_1 - x_0}x_1^2,\\
\Rightarrow\ \ \ 0 &= 2x_0^2x_1 - 2x_0x_1^2  - \frac{2}{3}(x_1-x_0),\\
\Rightarrow\ \ \ 0 &= 2x_0x_1 (x_1 - x_0) + \frac{2}{3}(x_1-x_0),
\end{align} 
also wegen $x_1\neq x_0$ dann $x_0x_1 = -\frac{1}{3}$.
Einsetzen in die letzte Gleichung liefert dann noch
\begin{align}
0 &= \frac{2x_1}{x_1 - x_0}x_0^3 - \frac{2x_0}{x_1 - x_0}x_1^3.
\end{align}
Die kompatible Lösung ist $x_1 = -x_0$, und damit $x_0 = -\frac{1}{\sqrt{3}}$ und $x_1 = \frac{1}{\sqrt{3}}$.
Damit sind dann $a_1 = 1 = a_0$.
Insgesamt lautet die Quadraturformel dann
\begin{align}
Q(f) &= f\kla\left( \frac{1}{\sqrt{3}} \right) + f\kla\left( -\frac{1}{\sqrt{3}} \right).
\end{align}
Die Formel funktioniert im Allgemeinen nicht mehr für Polynome höheren Grades.
Betrachten wir dazu einen zusätzlichen Term $+5x^4$, der im Integral dann den Beitrag
\begin{align}
I(f) &= \dots + x^5\Big\vert_{-1}^1 = \dots + 2
\end{align}
liefert.
Die Quadraturformel gibt allerdings
\begin{align}
Q(f) &= \dots + Q(5x^4) = \dots + \frac{10}{9},
\end{align}
was nicht mit dem Integralwert übereinstimmt.

\newpage
\section*{Aufgabe 4}
Sei $G = \{\alpha\in\mathbb{R}\, |\, \alpha > 0\}$ und 
\begin{align}
a &: G\times G \rightarrow G,\ a(x,y) = x+y,\\
m &: G\times G \rightarrow G,\ m(x,y) = xy.
\end{align}
\subsection*{a)}
Zu zeigen ist zunächst, dass diese Funktionen wohldefiniert sind. 
Insbesondere sollte geprüft werden, dass $a$ und $m$ jeweils tatsächlich nur auf Punkte in $G$ abbilden können.
Seien also $x,y\in G$ und $\alpha := a(x,y) = x + y$.
Dann gilt $x> 0$ und $y>0$, also auch $\alpha = x+y > 0$, und somit $\alpha\in G$.
Analog ist natürlich auch $\beta := m(x,y) = xy > 0$, also $\beta\in G$ (Anmerkung: Wir müssen die Aufgabe falsch verstanden haben, dieser ``Nachweis'' wird wohl kaum gefordert gewesen sein...).

\subsection*{b)}
Da $a$ und $m$ stetig differenzierbar sind, ist die absolute Kondition gegeben durch
\begin{align}
\kappa_\tecxt\text{abs}(a) &= \|a'(x,y)\| = \left\| \begin{pmatrix}
\frac{\partial a}{\partial x} & \frac{\partial a}{\partial y}
\end{pmatrix} \right\| = \left\|\begin{pmatrix}
 1 & 1
\end{pmatrix} \right\| = 1,
\end{align}
wobei wir die $\|\cdot\|_\infty$-Norm angenommen haben.
Analog gilt
\begin{align}
\kappa_\tecxt\text{abs}(m) &= \|m'(x,y)\| = \left\| \begin{pmatrix}
\frac{\partial m}{\partial x} & \frac{\partial m}{\partial y}
\end{pmatrix} \right\| = \left\|\begin{pmatrix}
 y & x
\end{pmatrix} \right\| = \max\{x,y\}.
\end{align}
Sei $v=(x,y)^\top\in G^2$. 
Für die relative Kondition gilt dann
\begin{align}
\kappa_\tecxt\text{rel}(a) &= \frac{\|v\|}{\|a(v)\|} \kappa_\tecxt\text{abs}(a) = \frac{\max\{x,y\}}{x+y},\\
\kappa_\tecxt\text{rel}(m) &= \frac{\|v\|}{\|m(v)\|} \kappa_\tecxt\text{abs}(m) = \frac{\max\{x,y\}}{xy} \max\{x,y\} = \klammer\left\{ \frac{x}{y}, \frac{y}{x} \right\}.
\end{align}



\newpage
\section*{Aufgabe 5}
Gegeben ist die Funktion $f:(0,\infty)\rightarrow\mathbb{R}$ mit 
\begin{align}
f(x) &= x + \log x.
\end{align}

\subsection*{a)}
Wir wollen zeigen, dass $f$ auf dem Intervall $I=(0,1)$ mindestens eine Nullstelle hat.
Da $f$ stetig auf $I$ ist, ist es hinreichend zu zeigen, dass $f$ auf $I$ einen Vorzeichenwechsel aufweist.
Für $x=\frac{1}{\e} \in I$ gilt offensichtlich 
\begin{align}
f(x) &= \frac{1}{\e} - 1 < 0,
\end{align}
da $\e > 1$. 
Für $x=1$ ist hingegen
\begin{align}
f(x) &= 1 > 0,
\end{align}
insbesondere gibt es also auch ein $x\in I$ mit $f(x) > 0$, da $f$ bei $x=1$ ansonsten einen Sprung hätte, und somit nicht stetig wäre.
Insbesondere ist
\begin{align}
f'(x) &= 1 + \frac{1}{x} > 0
\end{align}
für alle $x\in I$, die Funktion ist also streng monoton wachsend.
Damit kann es auf $I$ nur den einen Vorzeichenwechsel geben, den wir nachgewiseen haben.
Also hat $f$ auf $I$ genau eine Nullstelle $x^\ast$.

\subsection*{b)}
Um zu prüfen, ob $x^\ast$ ein Fixpunkt einer gegebenen Funktion $\phi(x)$ ist, können wir $x^\ast$ einfach einsetzen.
Dabei nutzen wir $0 = x^\ast + \log x^\ast$ aus.
Es gilt
\begin{align}
\phi_1(x^\ast) &:= -\log x^\ast = - (-x^\ast) = x^\ast,\\
\phi_2(x^\ast) &:= \e^{-x^\ast} = \e^{\log x^\ast} = x^\ast,\\
\phi_3(x^\ast) &:= x^\ast + \e^{-x^\ast} = 2x^\ast,\\
\phi_4(x^\ast) &:= \frac{1}{2}\kla\left( x^\ast + \e^{-x^\ast} \right) = x^\ast.
\end{align}
Also ist $x^\ast$ für $\phi_{1,2,4}$ tatsächlich ein Fixpunkt.
Um zu beweisen, dass es eine kleine Umgebung $I_r := [x^\ast - r, x^\ast + r]$ gibt, sodass die Fixpunktiteration $x_{k+1} := \phi (x_k)$ für beliebige Startwerte $x_0\in I_r$ gegen $x^\ast$ (bzw. überhaupt) konvergiert, können wir den Banach'schen Fixpunktsatz verwenden.
Voraussetzungen ist dabei die Lipschitz-Stetigkeit von $\phi$ auf $I_r$ mit Konstanter $L<1$.
Es gilt
\begin{align}
L_{\phi_2} &= \|\phi_2'(x)\| = \max_{x\in I_r} \e^{-x} < 1,\\
L_{\phi_4} &= \|\phi_4'(x)\| = \max_{x\in I_r} \frac{1}{2}\kla\left( 1 - \e^{-x} \right) < \frac{1}{2}.
\end{align}
Insbesondere gilt damit für $j\in\{2,4\}$ und $x\in I_r$
\begin{align}
|\phi_j(x) - x^\ast| &= |\phi_j(x) - \phi_j(x^\ast)| \le L_{\phi_j} |x-x^\ast| < |x-x^\ast| \le r,
\end{align}
also ist $\phi_j(x)$ weniger als $r$ entfernt von $x^\ast$, und damit immer noch in $I_r$.
Damit sind die Voraussetzungen des Fixpunktsatzes erfüllt.
\\
Für $\phi_3$ brauchen wir die Betrachtung gar nicht durchzuführen, da $x^\ast$ kein Fixpunkt ist. 
Für $\phi_1$ nehmen wir zunächst an, dass es ein beliebig kleines $r>0$ gebe, sodass die Folge $x_k$ gegen $x^\ast$ konvergiert.
Sei $x_0 = x^\ast + r$. 
Dann gilt
\begin{align}
\phi(x_0) &= -\log \kla\left( x^\ast + r \right) = -\log x^\ast + r (-\log x)'\Big\vert_{x=x^\ast} + \mathcal{O}(r^2) = \\
&= x^\ast - \frac{r}{x^\ast} + \mathcal{O}(r^2) =: x^\ast - r_1 + \mathcal{O}(r^2).
\end{align}
Insbesondere ist $r_1 = \frac{r}{x^\ast} > r$.
Der Fehler wächst also, selbst wenn man beliebig nahe an $x^\ast$ startet. 
Ein sauberer Beweis ist das nicht, aber mehr ist uns dazu nicht eingefallen.

\subsection*{c)}
Für die Iterationsvorschrift
\begin{align}
x_{k+1} &:= \frac{x_k^2 + \e^{-x_k}}{x_k + 1}
\end{align}
soll gezeigt werden, dass sie für alle $x\in I_r$ für ein hinreichend kleines $r>0$ quadratisch gegen $x^\ast$ konvergiert.
Dazu beobachtet man, dass die Newton-Vorschrift für die Funktion $g(x) := x\e^{x} - 1$ gerade durch
\begin{align}
x_{k+1} &= x_k - \frac{g(x_k)}{g'(x_k)} = x_k - \frac{x_k\e^{x_k} - 1}{\e^{x_k} + x_k\e^{x_k}} = x_k - \frac{x_k - \e^{-x_k}}{x_k + 1} = \frac{x_k^2 + \e^{-x_k}}{x_k + 1}
\end{align}
gegeben ist, was mit obiger Vorgabe übereinstimmt.
Da $g'(x) > 0$ für alle $x\in I$ ist, ist das Verfahren durchführbar und konvergiert damit quadratisch gegen die Nullstelle von $g$.
Diese stimmt mit $x^\ast$ überein, da 
\begin{align}
g(x) &= x\e^x - 1 \overset{!}{=} 0 &&\Rightarrow& x &= \log\frac{1}{x} && \Rightarrow & 0 &= x + \log x.
\end{align}

\newpage
\section*{Aufgabe 6}
Gegeben sind die Werte in folgender Tabelle. 
\begin{table}[!ht]
\centering
\begin{tabular}{l|c|c|c|c}
Zeit $t$ (in Stunden) & 0 & $0,5$ & 1 & $1,5$\\ \hline
Temperatur $T$ (in $^\circ$C) & 90 & 58 & 43 & 32\\
\end{tabular}
\end{table}
\\
Wir vermuten einen exponentiellen Zusammenhang der Form
\begin{align}
T(t) &= a\cdot 4^{-t} + b.
\end{align}
Das Modell ist linear in den zu bestimmenden Koeffizienten $x := (a,b)^\top$.
Wir können die Fehlerquadrate daher als
\begin{align}
\Delta^2 &= \| b - Ax\|^2
\end{align}
schreiben, wobei 
\begin{align}
b &= \begin{pmatrix}
T_1 \\ \vdots \\ T_4
\end{pmatrix} &&\tecxt\text{und}& A &= \begin{pmatrix}
4^{-t_1} & 1 \\ \vdots & \vdots \\ 4^{-t_4} & 1
\end{pmatrix}.
\end{align}
Dieser Ausdruck wird minimal für den Vektor $x^\ast\in\mathbb{R}^2$, der die Normalengleichung
\begin{align}
A^\top Ax^\ast &= A^\top b
\end{align}
löst.
%Die explizite Darstellung von $A$ und $A^\top A$ ist hier
%\begin{align}
%A &= \begin{pmatrix}
%1 & 1 \\ \frac{1}{2} & 1 \\ \frac{1}{4} & 1 \\ \frac{1}{8} & 1
%\end{pmatrix}\\
%A^\top A &= \begin{pmatrix}
%1 & 1 & 1 & 1 \\
%1 & \frac{1}{2} & \frac{1}{4} & \frac{1}{8}
%\end{pmatrix}\begin{pmatrix}
%1 & 1 \\ \frac{1}{2} & 1 \\ \frac{1}{4} & 1 \\ \frac{1}{8} & 1
%\end{pmatrix} = \begin{pmatrix}
%\frac{8+4+2+1}{8} & 4 \\ \frac{256 + 64 + 8 + 1}{256} & \frac{8+4+2+1}{8}
%\end{pmatrix} = \begin{pmatrix}
%\frac{15}{8} & 4 \\ \frac{329}{256} & \frac{15}{8}
%\end{pmatrix}.
%\end{align}
%Mithilfe der Pseudo-Inversen $A^+ := (A^\top A)^{-1}A^\top$, explizit
%\begin{align}
%A^+ &= \frac{1}{\kla\left( \frac{15}{8} \right)^2 - 4\cdot\frac{329}{256}} \begin{pmatrix}
%\frac{15}{8} & -4 \\ -\frac{329}{256} & \frac{15}{8}
%\end{pmatrix} A^\top = \\
%&= \frac{1}{\frac{225 - 329}{64}} \begin{pmatrix}
%\frac{15-32}{8} & \frac{15 - 16}{8} & \frac{15 - 8}{8} & \frac{15 - 4}{8} \\
%\frac{-329 + 15\cdot 32}{256} & \frac{-329 + 15\cdot 16}{256} & \frac{-329 + 15\cdot 8}{256} & \frac{-329 + 15\cdot 4}{256}
%\end{pmatrix} = \\
%&= \frac{64}{-104} \begin{pmatrix}
%\frac{-17}{8} & \frac{-1}{8} & \frac{7}{8} & \frac{11}{8} \\
%\frac{171}{256} & \frac{-89}{256} & \frac{-209}{256} & \frac{-269}{256}
%\end{pmatrix} = \frac{1}{416} \begin{pmatrix}
%544 & 32 & -224 & -352\\
%-171 & 89 & 209 & 269
%\end{pmatrix}.
%\end{align} 
Die explizite Darstellung von $A$ und $A^\top A$ ist hier
\begin{align}
A &= \begin{pmatrix}
1 & 1 \\ \frac{1}{2} & 1 \\ \frac{1}{4} & 1 \\ \frac{1}{8} & 1
\end{pmatrix}\\
A^\top A &= \begin{pmatrix}
1 & \frac{1}{2} & \frac{1}{4} & \frac{1}{8} \\
1 & 1 & 1 & 1 
\end{pmatrix}\begin{pmatrix}
1 & 1 \\ \frac{1}{2} & 1 \\ \frac{1}{4} & 1 \\ \frac{1}{8} & 1
\end{pmatrix} = \begin{pmatrix}
\frac{64+16+4+1}{64} & \frac{8+4+2+1}{8} \\ \frac{8+4+2+1}{8} & 4
\end{pmatrix} = \begin{pmatrix}
\frac{85}{64} & \frac{15}{8} \\ \frac{15}{8} & 4
\end{pmatrix}
\end{align}
Mithilfe der Pseudo-Inversen $A^+ := (A^\top A)^{-1}A^\top$, explizit
\begin{align}
A^+ &= \frac{1}{\frac{85}{16} - \kla\left( \frac{15}{8} \right)^2} \begin{pmatrix}
4 & -\frac{15}{8} \\ -\frac{15}{8} & \frac{85}{64}
\end{pmatrix} A^\top = \\
&= \frac{1}{\frac{340 - 225}{64}} \begin{pmatrix}
\frac{32 - 15}{8} & \frac{16 - 15}{8} & \frac{8 - 15}{8} & \frac{4 - 15}{8} \\
\frac{-15\cdot 8 + 85}{64} & \frac{-15\cdot 4 + 85}{64} & \frac{-15\cdot 2 + 85}{64} & \frac{-15\cdot 1 + 85}{64}
\end{pmatrix} = \\
&= \frac{64}{115} \begin{pmatrix}
\frac{17}{8} & \frac{1}{8} & -\frac{7}{8} & -\frac{11}{8} \\
\frac{-35}{64} & \frac{25}{64} & \frac{55}{64} & \frac{70}{64}
\end{pmatrix} = \frac{1}{115} \begin{pmatrix}
136 & 8 & - 56 & -88 \\
-35 & 25 & 55 & 70
\end{pmatrix}.
\end{align} 
lautet die Lösung dann
\begin{align}
x^\ast &= A^+ b = \frac{1}{115} \begin{pmatrix}
136 & 8 & - 56 & -88 \\
-35 & 25 & 55 & 70
\end{pmatrix} \begin{pmatrix}
90 \\ 58 \\ 43 \\ 32
\end{pmatrix} = \frac{1}{115} \begin{pmatrix}
7480 \\ 2905
\end{pmatrix}.
\end{align}
Die Auslgeichskurve lautet also
\begin{align}
T(t) &= \frac{1496}{23} \cdot 4^{-t} + \frac{581}{23}.
\end{align}
Nach 45 Minuten ($t= \frac{3}{4}$) lautet die extrapolierte Temperatur (in $^\circ$C) damit
\begin{align}
T(t) &= \frac{1496}{23} \cdot 2^{-3/2} + \frac{581}{23} \approx 48.26.
\end{align}
Abbildung \ref{figure:extrapolation} veranschaulicht nochmal die Extrapolation.
\begin{figure}[!ht]
\centering
\includegraphics{ExtrapolationWasser.png}
\caption{Visualisierung der Extrapolation für die gegebenen Werte. 
Der extrapolierte Wert der Wassertemperatur nach 45 Minuten ist hervorgehoben.}
\label{figure:extrapolation}
\end{figure}

%import matplotlib.pyplot as plt
%from matplotlib import special
%special.setup(UseTex=True)
%
%T = np.array([90, 58, 43, 32])
%t = np.arange(0, 2, 0.5)
%
%
%def model(t, a, b):
%    return a * 4.0**(-t) + b
%
%A = np.ones((4,2))
%A[:,0] = 4.0**(-t)
%Astar = np.linalg.inv(A.T @ A) @ A.T
%x = Astar @ T
%
%N = 300
%tval = np.linspace(0.0, 1.7, N)
%Tval = model(tval, *x)
%
%tc = 3/4
%Tc = model(tc, *x)
%
%fig, ax = plt.subplots()
%
%ax.set_xlabel(r"$t\ /\ $h")
%ax.set_ylabel(r"$T\ /\ ^\circ$C")
%ax.set_title(f"Extrapolation: $T(t) = {x[0]:.2f} \\cdot 4^{{-t}} + {x[1]:.2f}$")
%
%ax.set_xlim(tval[0], tval[-1])
%ax.set_ylim(0, np.max(b) * 1.05)
%
%ax.plot(t, b, ls='', marker='x', c='k', label="Messdaten")
%ax.plot(tval, Tval, c='b', label="Extrapolation")
%
%dt = 0.01 * (tval[-1] - tval[0])
%dT = 0.01 * (np.max(Tval) - np.min(Tval))
%ax.axvline(tc, c='k', ls='--', lw=0.7)
%ax.axhline(Tc, c='k', ls='--', lw=0.7)
%ax.text(tc + dt, Tc + 3.5*dT, f"$T\\left(t= \\frac{{3}}{{4}}\\right) \\approx {Tc:.2f}$")
%
%ax.legend(numpoints=1)
%special.polish(fig, ax)


%\begin{thebibliography}{9}
%\bibitem{example} description, \url{https://www.exmapleurl}, day.\,month~2021.
%\end{thebibliography}

\end{document}
